\chapter{Ibuprofen}

\section*{Overview}
Ibuprofen is a nonsteroidal anti-inflammatory drug (NSAID) used to treat pain, fever, and inflammation. It works by blocking the production of prostaglandins, chemicals that cause inflammation and pain in the body. Common brand names include Advil and Motrin.

\section*{Symptoms}
\begin{itemize}
  \item Pain relief (headaches, muscle aches, arthritis, menstrual cramps)
\end{itemize}

\begin{itemize}
  \item Fever reduction
\end{itemize}

\begin{itemize}
  \item Inflammation reduction
\end{itemize}

\begin{itemize}
  \item Toothache relief
\end{itemize}

\begin{itemize}
  \item Back pain
\end{itemize}

\section*{Causes}
Ibuprofen is not a condition but a medication. It is indicated for various conditions including osteoarthritis, rheumatoid arthritis, juvenile arthritis, mild to moderate pain, and fever.

\section*{Diagnosis}
Diagnosis is based on medical history, symptom assessment, and evaluation of pain severity. Healthcare providers assess the condition being treated rather than diagnosing with ibuprofen.

\section*{Treatment}
Take ibuprofen with food or milk to reduce stomach upset. Follow dosing instructions carefully. Do not exceed the recommended dosage. Adults typically take 200-400mg every 4-6 hours. Children's dosing is weight-based.

\section*{Prevention}
Prevent adverse effects by taking with food, staying hydrated, and avoiding alcohol. Use the lowest effective dose for the shortest duration necessary.

\section*{When to Seek Medical Care}
Seek medical attention if you experience stomach bleeding symptoms (black stools, vomiting blood), signs of allergic reaction (rash, swelling, difficulty breathing), or symptoms of kidney problems.
\clearpage
